\section{Introduction}

\subsection*{The reduction}
For an admissible $k$-tuple $\mathcal H=\{h_1<\dots<h_k\}$, let $\DHL[k,2]$
denote the statement that infinitely many translates $n+\mathcal H$ contain at
least two primes. If $\DHL[k,2]$ holds and $\mathcal H$ is an admissible
$k$-tuple of diameter $D$, then $H_1\le D$. Maynard's method \cite{May15}
proves $\DHL[k,2]$ from a level of distribution and a symmetric test function
$F$ whose ratio $kJ(F)/I(F)$ exceeds a threshold; Polymath8b \cite{Pol14b}
enlarged the region on which $F$ may live and obtained $H_1\le246$ with
$k=50$. The 2026 preprints listed below enlarge the region further, to
\emph{capped} supports, and pay for this with equidistribution estimates
for the moduli that arise: Stadlmann \cite{Sta26} proves new ones, Kintali
\cite{Kin26} uses Stadlmann's \cite[Prop.~3]{Sta26} with one branch made empty,
Charton et al.\ \cite{Cha26} use Stadlmann's Type~I, II and III estimates for
their own family of moduli, and Song and Yue \cite{SY26} prove their own.

\subsection*{The claims examined}
\begin{center}
\begin{tabular}{@{}llll@{}}
\toprule
claim & $k$ & source & what we check\\
\midrule
$H_1\le240$ & $49$ & Stadlmann \cite{Sta26} & variational criterion\\
$H_1\le236$ & $48$ & Kintali \cite{Kin26}; Song--Yue \cite{SY26} & variational criterion; published witness\\
$H_1\le212$ & $45$ & Charton et al.\ \cite{Cha26} & variational criterion; caps\\
\bottomrule
\end{tabular}
\end{center}
Each claim has two parts: an \emph{analytic} part (an equidistribution
theorem for the relevant family of moduli) and a \emph{variational} part (a
function $F$ on the capped support with ratio above the threshold). We check
only the variational part, and the combinatorial fact that an admissible
tuple of the stated diameter exists. Throughout we label statements
\lbl{Computed} (exact certificate and a script that reproduces it),
\lbl{Proved here}, \lbl{Lean-checked}, \lbl{Literature} (read in the source),
or \lbl{Claimed} (asserted in a preprint and not verified by us).

\subsection*{Results in brief}
Ratios in the Polymath scaling, where the threshold is $4$:
\begin{center}\small
\begin{tabular}{@{}lll@{}}
\toprule
claim & variational ratio & status here\\
\midrule
$240$ & $4.0120389$; optimised on $T$: $4.0124198$ & \lbl{Computed}, two programs\\
$236$, witness of \cite{SY26} & $48J_T/I_T=4.0030988$; $48J_T/I_0=4.0026597$ & \lbl{Computed}, two programs\\
 & \quad($=4\times1.00066492$; SY criterion $>2001/2000$) & ($J_T$; $I_T$ one)\\
$236$, our own $F$ & $4.0060814$ & \lbl{Computed}, two programs\\
$212$, Table~3 caps & $4.0027351$ & \lbl{Computed}, two programs\\
$212$, Lemma~9.1 caps & $<3.99$ for all $F$ in an $846$-dim.\ space & \lbl{Computed}, ball arithm.\\
\bottomrule
\end{tabular}
\end{center}

\subsection*{What this note does not do}
It does not verify the equidistribution estimates of \cite{Sta26},
\cite{Kin26}, \cite{SY26} or \cite{Cha26}; all three claimed bounds remain conditional
on them, and the established bound is still $246$. Song and Yue
\cite[Remark~4.23]{SY26} point to a missing hypothesis in one step of the
proof of \cite[Lemma~5]{Sta26}; we read that step and comment on it, and on
which of the claims may depend on it, in Remark~\ref{rem:lemma5}, without
resolving it. Nothing here concerns the
parity barrier or the twin prime conjecture.

